\section*{Sets and Indices}

Let
\[
T = \{\text{robot},\text{car},\text{blocks},\text{doll}\}
\]
be the set of toy types.

No additional indexing sets are used in the implemented model.

\section*{Parameters}

Profits per unit:
\[
p_{\text{robot}}=15,\quad
p_{\text{car}}=8,\quad
p_{\text{blocks}}=12,\quad
p_{\text{doll}}=5.
\]

Plastic consumption per unit:
\[
a_{\text{robot}}=30,\quad
a_{\text{car}}=10,\quad
a_{\text{blocks}}=20,\quad
a_{\text{doll}}=15.
\]

Available plastic:
\[
A=1200.
\]

Big-\(M\) constant used for logical linking:
\[
M=10000.
\]

The CSV also contains electronic-component data and Boolean flags for business rules, but in the current implementation only the plastic parameters and the three logical rules listed below are enforced.

\section*{Decision Variables}

Integer production quantities:
\[
x_{\text{robot}} \in \mathbb{Z}_{\ge 0} \quad \text{(code: \texttt{x\_robot})},
\]
\[
x_{\text{car}} \in \mathbb{Z}_{\ge 0} \quad \text{(code: \texttt{x\_car})},
\]
\[
x_{\text{blocks}} \in \mathbb{Z}_{\ge 0} \quad \text{(code: \texttt{x\_blocks})},
\]
\[
x_{\text{doll}} \in \mathbb{Z}_{\ge 0} \quad \text{(code: \texttt{x\_doll})}.
\]

Binary indicators for whether a toy type is manufactured:
\[
y_{\text{robot}} \in \{0,1\} \quad \text{(code: \texttt{y\_robot})},
\]
\[
y_{\text{car}} \in \{0,1\} \quad \text{(code: \texttt{y\_car})},
\]
\[
y_{\text{blocks}} \in \{0,1\} \quad \text{(code: \texttt{y\_blocks})},
\]
\[
y_{\text{doll}} \in \{0,1\} \quad \text{(code: \texttt{y\_doll})}.
\]

\section*{Objective}

Maximize total profit:
\[
\max \;
p_{\text{robot}}x_{\text{robot}}
+ p_{\text{car}}x_{\text{car}}
+ p_{\text{blocks}}x_{\text{blocks}}
+ p_{\text{doll}}x_{\text{doll}}.
\]

With the given data:
\[
\max \;
15x_{\text{robot}} + 8x_{\text{car}} + 12x_{\text{blocks}} + 5x_{\text{doll}}.
\]

\section*{Constraints}

Plastic capacity:
\[
a_{\text{robot}}x_{\text{robot}}
+ a_{\text{car}}x_{\text{car}}
+ a_{\text{blocks}}x_{\text{blocks}}
+ a_{\text{doll}}x_{\text{doll}}
\le A.
\]

With the given data:
\[
30x_{\text{robot}} + 10x_{\text{car}} + 20x_{\text{blocks}} + 15x_{\text{doll}} \le 1200.
\]

Upper linking of production to activity indicators:
\[
x_{\text{robot}} \le M y_{\text{robot}},
\]
\[
x_{\text{car}} \le M y_{\text{car}},
\]
\[
x_{\text{blocks}} \le M y_{\text{blocks}},
\]
\[
x_{\text{doll}} \le M y_{\text{doll}}.
\]

Lower linking of activity indicators to production:
\[
y_{\text{robot}} \le x_{\text{robot}},
\]
\[
y_{\text{car}} \le x_{\text{car}},
\]
\[
y_{\text{blocks}} \le x_{\text{blocks}},
\]
\[
y_{\text{doll}} \le x_{\text{doll}}.
\]

Robot--doll mutual exclusion:
\[
y_{\text{robot}} + y_{\text{doll}} \le 1.
\]

Model-car implies building-blocks production:
\[
y_{\text{car}} \le y_{\text{blocks}}.
\]

Dolls cannot exceed model cars:
\[
x_{\text{doll}} \le x_{\text{car}}.
\]

Integrality and binary restrictions:
\[
x_t \in \mathbb{Z}_{\ge 0} \qquad \forall t \in T,
\]
\[
y_t \in \{0,1\} \qquad \forall t \in T.
\]

\section*{Notes on Implementation}

The implemented model is a mixed-integer linear program solved via \texttt{GUROBI\_CMD} through a PuLP-compatible interface.

The formulation intentionally matches the code exactly:
\begin{itemize}
\item Only the plastic resource constraint is enforced.
\item Although the data file provides electronic-component capacities and requirements, they are not used in the current model.
\item The three business rules are modeled through binary variables and linear constraints exactly as coded.
\item The pair of constraints \(x_t \le M y_t\) and \(y_t \le x_t\) forces \(y_t=1\) iff \(x_t \ge 1\), since \(x_t\) is a nonnegative integer and \(y_t\) is binary.
\item The big-\(M\) value is fixed in code as \(M=10000\); it is not derived from resource bounds.
\end{itemize}
